% 原文第37页；损耗与补偿公式、时序图保留。
\zhbox{54}{104}{235}{26}{9.5}{SWA 的开关损耗 $P_{SW\_LOSS\_SWA}$ 可由式（19）计算：}
\zhbox{54}{170}{235}{26}{9.5}{MOSFET 的开通时间 $t_{ON}$ 和关断时间 $t_{OFF}$ 参见图13。}
\zhbox{57}{333}{232}{13}{9.5}{\centering\textbf{图13：开通时间与关断时间}}
\zhbox{54}{351}{235}{54}{9.5}{\textbf{SWB MOSFET}\\降压模式下，SWB 的导通损耗 $P_{CON\_LOSS\_SWB}$ 可由式（20）计算：}
\zhbox{54}{450}{235}{55}{9.5}{\textbf{SWC MOSFET}\\升压模式下，SWB 关断，SWC 的导通损耗 $P_{CON\_LOSS\_SWC}$ 可由式（21）计算：}
\zhbox{54}{545}{235}{38}{9.5}{升压模式下，SWC 的开关损耗 $P_{SW\_LOSS\_SWC}$ 可由式（22）计算：}
\zhbox{54}{620}{235}{54}{9.5}{\textbf{SWD MOSFET}\\降压模式下，SWD 导通，其功率损耗 $P_{CON\_LOSS\_SWD}$ 可由式（23）计算：}
\zhbox{324}{64}{235}{26}{9.5}{升压模式下，SWD 的导通损耗可由式（24）计算：}
\zhbox{324}{128}{235}{26}{9.5}{可忽略死区时间 $t_{DEAD}$ 和低侧 MOSFET 的开关损耗。}
\zhbox{324}{159}{235}{79}{9.5}{\textbf{补偿元件}\\COMP 引脚控制系统稳定性和瞬态响应。COMP 为内部误差放大器（EA）的输出。串联的电容和电阻设置一个极点和一个零点，以控制系统特性。}
\zhbox{340}{348}{210}{14}{9.5}{\centering\textbf{图14：COMP 外部补偿}}
\zhbox{324}{366}{235}{39}{9.5}{外部补偿设置一个极点 $F_{P1}$ 和一个零点 $F_{Z1}$（见图14）。$F_{P1}$ 可由式（25）计算：}
\zhbox{324}{447}{235}{14}{9.5}{$F_{Z1}$ 可由式（26）计算：}
\zhbox{324}{502}{235}{27}{9.5}{降压模式下，电压反馈环路的直流增益 $A_{VDC}$ 可由式（27）计算：}
\zhbox{324}{570}{235}{52}{9.5}{其中，$G_{V\_EA}$ 为 EA 电压增益（300V/V），$G_{CS}$ 为 COMP 到电流检测的增益，$R_{SENSE1}$ 为电流检测电阻，$R_{LOAD}$ 为负载电阻。}
\zhbox{324}{627}{235}{78}{9.5}{系统有两个重要极点：一个由补偿电容 C10 和 EA 输出电阻产生，另一个由输出电容 $C_{OUT}$ 和负载电阻产生。补偿电容极点 $F_{P2}$ 可由式（28）计算：}
